This study evaluated whether selected historical states of Bangladesh's public DGHS Dengue Dynamic Dashboard could be reconstructed, archived, and interpreted as research-data artifacts. Within the validated corpus, the workflow supported traceable historical-state reconstruction, SHA-256-linked provenance, source-labelled extraction, and explicit recording of category and reconciliation limitations. It also demonstrated that reproducible retrieval does not by itself resolve case or admission semantics, reporting-facility coverage, geographic attribution, denominator compatibility, incomplete 2026 follow-up, or retrospective revision behavior.

The practical implication is that public epidemiological dashboards can support bounded secondary research when their observation universe is stated clearly and their raw responses, retrieval context, labels, and limitations are preserved. Researchers should validate semantics before normalization, compare fields only when their time windows and reporting universes are compatible, and distinguish dashboard observations from the underlying surveillance records that generated them. In this setting, provenance, semantic validation, reconciliation, and documentation of historical-state changes are prerequisites for defensible reuse, not optional post-processing steps.
